% Chapter 2: Related Work
% Citations use natbib (\citet, \citep). Bibliography entries: thesis/references_ch2.bib
% Body text target ~1,700 words of the 8,000-word limit (appendix excluded).

\chapter{Related Work}
\label{ch:related}

This chapter covers the coding schemes and how reliably people apply them, automatic coders from n-gram models to large language models (LLMs), and the method families this thesis tests: local parameter-efficient adaptation, rationale and reinforcement-learning training, and ways to add labels without another round of expert coding. Section~\ref{sec:rw-gaps} maps the open problems onto the research questions.

\section{Coding motivational interviewing}
\label{sec:rw-coding}

Motivational interviewing (MI) is a counselling style that evokes a client's own reasons for change \citep{miller2023mi}. Across 119 studies it produced small, durable effects against weak comparison groups \citep{lundahl2010meta}. Fidelity coding matters because of the technical hypothesis: specific counsellor behaviours shape client language, and client language predicts outcome. \citet{magill2014technical} gave the first aggregate test. In a later meta-analysis of 36 studies, MI-consistent skills correlated with more change talk, a higher proportion of change talk correlated with less risk behaviour at follow-up, and sustain talk predicted worse outcome \citep{magill2018process}.

Some of the behaviours that carry this effect are rare in a session. Affirm was the only counsellor behaviour that both raised change talk and lowered sustain talk \citep{apodaca2016which}. An earlier review placed MI-inconsistent behaviour, such as confronting or advising without permission, among the few constructs consistently linked to worse outcome \citep{apodaca2009mechanisms}. A coder that gets questions and simple reflections right but misses these codes gives feedback on the easy part of a session.

Two schemes dominate. The Motivational Interviewing Skill Code (MISC) 2.5 codes every utterance of both speakers and groups counsellor codes into MI-consistent and MI-inconsistent summaries \citep{houck2010misc}. The Motivational Interviewing Treatment Integrity code (MITI) was derived from MISC by factor analysis, and codes the counsellor only, with behaviour counts and global ratings \citep{moyers2005miti}; version~4 revised its codes and globals \citep{moyers2016miti4}. Public corpora use further variants. AnnoMI labels 133 transcribed demonstrations with coarse MISC-derived attributes \citep{wu2022annomi,wu2023annomi}; \citet{welivita2022curating} adapted MITI to about 17,000 peer-support forum responses; MI-TAGS codes 242 demonstrations with MITI~4.2 \citep{cohen2024mitags}; BiMISC adds real sessions in Dutch and English \citep{sun2024bimisc}. The schemes differ in granularity and in whether client speech is coded, so results transfer between them only through a code mapping. AutoMISC grouped MISC~2.5 codes into a two-tier hierarchy, a coarse group (T1) and a fine code (T2) \citep{ali2025automisc}. That grouping is AutoMISC's, not the manual's.

Human agreement sets the ceiling for any coder. In an early MISC study, 75\% of global ratings but only 44\% of behaviour counts reached good-to-excellent intraclass correlation \citep{moyers2003misc}. On a 277-session MITI corpus, session-level ICC ranged from 0.75 to 1, while utterance-level Cohen's $\kappa$ ranged from 0.31 to 0.64 \citep{perezrosas2016building}. \citet{lord2015advancing} found the same gap and argued that reliability estimated from session tallies overstates utterance-level agreement. Even MITI~4 fell short on Emphasize Autonomy and percent complex reflections \citep{moyers2016miti4}. Coding is also slow; MISC has been called labour-intensive \citep{dejonge2005misc}, and feedback that depends on trained coding teams is ``slow, expensive, and labor intensive'' \citep{imel2019design}. Two consequences follow. Utterance-level gold carries the conventions of the team that coded it, and session-level summaries are the more reliable target for feedback.

\section{Automatic coders before large language models}
\label{sec:rw-classic}

A labelled topic model reproduced MISC codes with AUC between 0.62 and 0.81 \citep{atkins2014scaling}, and a sequence model detected reflections with 93\% recall and 73\% precision \citep{can2016sounds}. On 341 sessions, \citet{tanana2016comparison} reached $\kappa > 0.60$ for open and closed questions, affirm, giving information and follow/neutral, but agreement was poor for change talk, sustain talk and MI-inconsistent behaviour. Neural models followed: feature-rich counsellor classifiers \citep{perezrosas2017predicting}, dialogue models that categorise and forecast codes \citep{cao2019observing}, and RoBERTa client classifiers with a three-class F1 of 0.66 \citep{tavabi2021analysis}. BERTje on 253 helpline chats reached $\kappa = 0.69$ for counsellor behaviour, with the caveat that each code needs enough training examples \citep{pellemans2024automated}. In school coaching transcripts the codes of most interest were the rarest and were coded with poor accuracy \citep{pas2025leveraging}.

The pattern has held for over a decade: frequent codes are learned, rare ones are not. It also matters outside benchmarks, because these coders feed real feedback systems \citep{imel2019design,flemotomos2022automated}. Transcript quality is a second source of error: speech-recognition errors affect counsellor coding, and extra local context reduces the damage \citep{min2021evaluating}.

\section{Large language models as coders}
\label{sec:rw-llm}

Prompted LLMs removed the need for task-specific training. \citet{sun2024bimisc} elicited MISC codes from an LLM that aligned closely with expert annotation, and Chain-of-Interaction prompting improved client coding by reasoning over the counsellor--client exchange \citep{han2024chain}. AutoMISC segments each turn into utterances and codes T1 then T2 with GPT-4.1 and three prior volleys of context; it reports counsellor T2 accuracy of 0.68 and macro-F1 of 0.42 \citep{ali2025automisc}. The 26-point gap between those two numbers is the long tail again. AutoMISC's human-coded sets are also the MISC training and test data of this thesis (Chapter~4). A GPT-4.1 coder in a prevention programme agreed with humans on affirmations ($\kappa$ 0.66--0.77) but not on complex reflections ($\kappa = 0.00$), although only two practice sessions were coded \citep{cardozo2026automating}.

For routine fidelity feedback the prompted route has two costs. Session text leaves the service for a third-party API, and the vendor can change the model, so a result cannot be re-run on fixed weights. Small open models avoid both and have been proposed as privacy-preserving tools for mental-health text \citep{jia2025beyond}. In general text classification, fine-tuned smaller models still beat zero-shot GPT-3.5, GPT-4 and Claude Opus \citep{bucher2024finetuned}, fine-tuning smaller models beats few-shot prompting across 16 datasets \citep{edwards2024language}, and fine-tuned open models match or surpass zero-shot GPT-4 on annotation tasks \citep{alizadeh2025opensource}. Within MI, instruction fine-tuning was more stable to prompt choice than in-context learning and improved with more data for client motivational language \citep{hoang2024client}, and fine-tuned models have labelled 17 MI techniques in peer-counselling chats \citep{shah2022modeling}. I did not find a study that fine-tunes an open LLM on expert MISC labels for the full two-tier code set of both speakers and compares it with a frontier prompt on the same held-out sessions.

\section{Local adaptation}
\label{sec:rw-peft}

Low-rank adaptation (LoRA) freezes the pretrained weights and trains low-rank update matrices in each layer. On GPT-3 it cut trainable parameters by 10,000 times and GPU memory by three times, matched full fine-tuning, and added no inference latency \citep{hu2022lora}. With instruction tuning, a classification task becomes text generation from a prompt \citep{wei2022finetuned}, so a 7--15B model can learn the code vocabulary on one GPU.

The two-tier hierarchy raises two design questions. The first is whether both tiers share one adapter. Training one task after another causes forgetting in LLMs \citep{luo2025empirical}, and LoRA does not prevent it; forgetting grows with the number of update steps \citep{kalajdzievski2024scaling}. The second is what T2 is conditioned on. If T2 is trained on gold T1 but sees predicted T1 at test time, the model meets inputs it never saw in training, the exposure-bias problem. Scheduled sampling mixes gold and predicted inputs during training \citep{bengio2015scheduled}, but it can make a model worse when the prefix at inference is correct \citep{korakakis2022improving}. In a two-call coder most T1 predictions are correct, so this trade-off applies directly.

\section{Rationales and reinforcement learning}
\label{sec:rw-reasoning}

Asking a model to reason before answering improves many tasks \citep{wei2022chain,kojima2022large}, and rationales from a large teacher can train a small student. With rationales as extra supervision, a 770M T5 model beat few-shot 540B PaLM using 80\% of the data on one benchmark \citep{hsieh2023distilling}; Fine-tune-CoT showed similar gains \citep{ho2023large}. STaR adds rationalisation: when the model answers wrongly, it writes a rationale with the correct answer given as a hint \citep{zelikman2022star}. Rationales written with the gold code visible, as in this thesis, are of that post-hoc kind.

A rationale is only a checkable reason if the model uses it. Chain-of-thought explanations can misstate the real cause of an answer, for example by omitting a biasing feature in the prompt \citep{turpin2023language}. When chains are corrupted or paraphrased, models differ widely by task in how much the answer depends on them \citep{lanham2023measuring}, and such tests measure output self-consistency rather than faithfulness to internal computation \citep{parcalabescu2024measuring}. For a coder, then, accuracy alone does not show that a trained rationale is load-bearing. The rationale has to be perturbed and the code watched, against a model never trained on rationales.

Group Relative Policy Optimization (GRPO) drops PPO's value model and scores each sample against the other samples drawn for the same prompt \citep{shao2024deepseekmath}. With verifiable rewards it produced strong reasoning models \citep{deepseek2025r1}. A gold code is a verifiable reward, which makes coding a natural fit. Two findings temper this. GRPO's loss inflates response length, which Dr.~GRPO corrects \citep{liu2025understanding}. And GRPO reinforces trajectories the model already finds likely while neglecting rare correct ones \citep{he2025rewarding}. For a long-tailed label space, that rank bias suggests sharper head-code predictions rather than recovered tail codes. Published GRPO results come mostly from mathematics and code, where answers are exact; coding against noisy, subjective gold is largely untested.

\section{Adding labels without another expert pass}
\label{sec:rw-labels}

Class re-weighting \citep{cui2019class} changes how much each rare example counts but adds no new examples. Two routes add examples. Self-training labels an unlabelled pool with the current model. Naive pseudo-labelling overfits to its own errors \citep{arazo2020pseudo}, and under class imbalance pseudo-labels drift toward the majority classes \citep{he2021redistributing}. CReST found minority pseudo-labels still precise and sampled them more often \citep{wei2021crest}, the same logic as leaving rare codes uncapped when selecting pseudo-labels (Chapter~3).

Synthesis asks an LLM to write labelled examples. Its value falls as the task becomes more subjective, at both task and instance level \citep{li2023synthetic}. Simple class-conditional prompts produce biased, low-diversity data \citep{yu2023large}, and pushing diversity up lowers label accuracy unless a human corrects the labels \citep{chung2023increasing}. MI coding is subjective and long-tailed, which is the hard case for both routes. IC-AnnoMI rewrote AnnoMI dialogues with ChatGPT and had an expert rate them at dialogue level, while the utterances kept their source labels \citep{kumar2024unlocking}. Without an expert check on each label, a model has to verify the labels itself.

A third, lighter option retrieves labelled examples into the prompt. Similar examples beat random ones \citep{liu2022makes}, but randomising the labels in demonstrations barely hurts accuracy; the label space, input distribution and format do most of the work \citep{min2022rethinking}. Retrieval may therefore teach the output format more than the boundary between two rare codes.

\section{Gaps and research questions}
\label{sec:rw-gaps}

Four gaps follow. First, prompted frontier models are the reference for MISC coding, but evidence for a locally adapted open model on the full two-tier code set is missing (RQ1, access). Second, every generation of coder, from n-gram models to GPT-4.1, does worst on rare codes, and accuracy hides it; whether local adaptation, or a change to the pipeline structure, helps the tail has to be measured per code (RQ2). Third, rationale distillation and GRPO were shown on tasks with exact answers, and faithfulness work shows a rationale may not be used. Reasoning prompts have been tried for MI coding \citep{han2024chain}, but I found no study that trains a coder on rationales or with GRPO and then tests whether its rationale is used (RQ3). Fourth, self-training drifts toward the majority and synthetic data weakens as subjectivity rises. MI coding is both subjective and long-tailed, and I found no test of either route on rare MISC codes against expert gold (RQ4).

The research proposal also planned active learning and a coaching corpus. Neither was run, and neither is evaluated here.
% TODO(Jiawen): confirm this list against the proposal PDF; add joint MISC--MITI training here if the re-run does not cover it.
